%!TEX root = ../template.tex
%%%%%%%%%%%%%%%%%%%%%%%%%%%%%%%%%%%%%%%%%%%%%%%%%%%%%%%%%%%%%%%%%%%%
%% chapter5.tex
%% NOVA thesis document file
%%
%% Conclusions and future work
%%%%%%%%%%%%%%%%%%%%%%%%%%%%%%%%%%%%%%%%%%%%%%%%%%%%%%%%%%%%%%%%%%%%

\typeout{NT FILE chapter5.tex}%

\chapter{Conclusions and Future Work}
\label{cha:conclusions}

This thesis asked whether the apparent robustness of MADDPG routing controllers to FGSM
survives an observation-space attack that coordinates its perturbation across agents,
inspired by SAJA, and divides it across time, inspired by PGD and by the
strategically-timed attack of Lin et al.~\cite{lin2019tacticsadversarialattackdeep}. The
attack was realised
as a learned adversary, trained against the frozen victim and evaluated against seven
victim architectures in paired episodes with a random control. This chapter draws the
conclusions of that evaluation on the four questions of
Section~\ref{sec:intro_problem}: how effective the attack is compared with the FGSM
baseline, which is the central question; what coordination adds over the independent
attack (RQ2); what the timing budget achieves compared with attacking every step (RQ1);
and which architectural choices of the victim determine how vulnerable it is (RQ3). It
then sets out
the limitations of those conclusions, above all their statistical relevance, and
outlines the work that would strengthen or extend them.

% ─────────────────────────────────────────────────────────────────────────────
\section{Conclusions}
\label{sec:conclusions}

\subsection{Effectiveness compared with the FGSM baseline}
\label{subsec:conclusions_fgsm}

The comparison with FGSM is only meaningful under the FGSM study's own admissible set,
and it was made on the 15 episodes that the two evaluations share
(Section~\ref{sec:results_fgsm_comparison}).

The learned attack did more total damage than FGSM in one configuration only, the
independent attack on every step against \texttt{LC-Simple}, by $2.21$~pp; everywhere
else it matched the baseline or fell short of it. With the timing budget in place, the
independent attack did significantly less damage than FGSM on four variants at every
budget, which is partly by construction, since it acted on only $10$ to $25\%$ of the
steps while FGSM acted on all of them. Coordination brought the attack closer to FGSM
on the duelling variants: on \texttt{LC-Duelling} it could not be distinguished from
FGSM at any budget, and on \texttt{CC-Duelling} from $b = 0.20$ onwards. When attacking
every step, the coordinated attack did not differ significantly from FGSM on any
variant. Where the learned attack has an advantage, it is per attacked step: on
\texttt{LC-Simple} in the independent scope, and on \texttt{CC-Duelling} in the
coordinated scope at $b = 0.25$, it caused more damage per step attacked than FGSM,
without causing more damage in total.

The victim family withstands the learned
adversary at least as well as it withstands the single-step baseline. The learned adversary
is theoretically the stronger attacker: its policy class includes FGSM's per-step rule,
and it can additionally coordinate across agents and choose when to act. That it does no
more damage than FGSM is consistent with a robust controller, and it is the reading this
thesis favours over the alternative, that the FGSM baseline was simply too weak. It does
not, however, establish robustness, for two reasons that run through the evaluation.

The first is the size of the effects. Against a clean delivery ratio of about $77\%$, no
gated gap in any configuration reaches $1.5$~pp, which is the same order as the movement
that a perturbation of this size produces with no adversarial direction at all: on
\texttt{LC-Simple} the random control delivered $0.4$~pp \emph{more} than the clean arm,
and on \texttt{LC-Duelling} the trained adversary delivered $1.7$~pp more than the clean
arm (Section~\ref{subsec:results_parity_timing}). The second is that each configuration
was trained once and without a seed, so the confidence intervals of
Chapter~\ref{chap:results} cover the variation between evaluation episodes but not the
variation between training runs (Section~\ref{subsec:limitations_statistics}). An effect
that clears the first hurdle but not the second says little about the controller.

What the evaluation establishes is therefore the weaker statement that no attack in this
family, as trained here, did substantial damage to these controllers. How much of the
variation it reports is the adversary's doing rather than noise, and whether controllers
of this kind are robust in general, remain open. A learned adversary is also only as
strong as its training, and one of them steered \texttt{LC-Duelling} to deliver more
packets than it did unattacked (Section~\ref{subsec:results_parity_timing}), which shows
that a trained adversary can fail to find a damaging perturbation where one exists.

\subsection{Coordination compared with the independent attack}
\label{subsec:conclusions_coordination}

Coordination changes the outcome of the attack, but it does not improve it uniformly.
Under the FGSM study's admissible set it changed the damage significantly on four
victims, in both directions (Section~\ref{subsec:results_parity_coordination_effect}).
It increased the damage done to \texttt{LC-Duelling} at every budget, by $0.64$ to
$1.58$~pp, turning a victim the independent adversary had helped into one the attack
harms, and to \texttt{CC-Duelling} from $b = 0.15$ onwards, by up to $2.16$~pp. It
reduced the damage done to \texttt{LC-Simple} at every budget, by $0.48$ to $3.99$~pp,
on the variant the independent attack had harmed most. Under the stricter full domain
constraint, the full clamp of Chapter~\ref{chap:results}, the
pattern was different again: coordination made \texttt{LC-Simple} attackable and
reduced the damage done to \texttt{CC-Duelling-GNN}
(Section~\ref{subsec:results_coordination_effect}).

The idea taken from SAJA, judging the perturbations of several agents together by
their combined effect, is therefore able to reach damage that the independent attack
does not, on some victims, while losing damage the independent attack had found on
others. No single property of the victim explains which: the local-critic account that
fits the full-clamp results does not fit the FGSM-parity ones. What coordination
contributes depends on the pair of trained adversaries being compared as much as on the
victim, and its gains do not come from changing more routing decisions.

That last point should be pressed further. Every coordination effect reported here is a
difference between two separately trained adversaries, and almost all of them are
smaller than $1.5$~pp. Since neither adversary's training was repeated, the spread
between training runs is unknown and cannot be subtracted from these differences
(Section~\ref{subsec:limitations_statistics}). The direction of the effect on a given
victim is therefore not established; what the evaluation does show is that coordination
is not a uniform improvement, and that no configuration of it produced a large gap.

\subsection{Timing budget compared with attacking every step}
\label{subsec:conclusions_timing}

The timing gate itself works as designed. It realised the requested budgets closely,
at $10.0\%$, $15.0\%$, $19.9\%$ and $24.9\%$ of steps on average, and the trigger
analysis of Sections~\ref{subsec:results_parity_triggers}
and~\ref{subsec:results_parity_coord} showed why the tie-rationing rule is needed. The
critical-state score repeatedly takes the same value of $0.90$, and without rationing
the two tightest budgets would both have attacked between $17\%$ and $21\%$ of steps.
In the FGSM-parity runs, where the reasons were recorded, the attacks fired only when the
most utilised link was at least $75\%$ utilised, and at the two tightest budgets almost
entirely at $90\%$ or more. The composition of the
triggers matched to within a few percentage points in both scopes, so it is set by the
network and the budget rather than by the perturbation.

Spending a limited budget on those moments did not buy more total damage than
attacking every step. For almost every victim the attack harmed, the ungated attack did
the most damage. What the budget can buy is efficiency: in three of the four settings
evaluated, the damage per attacked step was higher under a tight budget than when
attacking every step, by a factor of up to $2.6$ on \texttt{CC-Duelling-GNN} under the
full clamp and of about $1.3$ on \texttt{CC-Duelling} in the coordinated scope under
the FGSM study's admissible set. It was not higher for the independent attack under that
admissible set, where \texttt{LC-Simple}'s damage grew in proportion to the number of
steps attacked. The idea taken from PGD and from Lin et al., of dividing the attack
across time and spending it at chosen moments, therefore
makes each attacked step count for more in most settings, but not reliably, and it does
not make the attack more damaging overall. These efficiency figures divide one small
quantity by another, a gap of at most a few percentage points by an attack rate of a
tenth to a quarter, and carry no confidence interval, so they are indications rather
than measurements, and the factors quoted above should not be read as precise. A further
limitation of the ungated comparison is that every adversary was trained with the gate at
$b = 0.25$, so when attacking every step it acted on states it had not encountered during
training.

\subsection{Architectural choices and vulnerability}
\label{subsec:conclusions_architecture}

No architectural axis of the victim family is vulnerable across the board. Which choice
appears to be the weak one changes with the admissible set and with the attack scope.
Under the FGSM study's set, the independent attack harmed \texttt{LC-Simple} at every
budget and \texttt{CC-Duelling-GNN} only when attacking every step, while the
coordinated attack harmed \texttt{CC-Duelling} most
(Sections~\ref{subsec:results_parity_timing} and~\ref{subsec:results_parity_coord}).
Under the full clamp the pattern moved again: the independent attack harmed the duelling
variants, led by \texttt{CC-Duelling-GNN} and \texttt{LC-Duelling}, and the coordinated
attack harmed the two local-critic variants without GNN encoding
(Sections~\ref{subsec:results_ind_timing} and~\ref{subsec:results_coord}). Each of the
three axes therefore appears decisive in some setting and irrelevant in another, so none
of them can be named as the vulnerability of the controller, and a defence chosen on the
evidence of one setting would be chosen on the wrong evidence. The caution of
Section~\ref{subsec:conclusions_fgsm} applies here with particular force: these are
comparisons between separately trained adversaries at gaps of one or two percentage
points, so the pattern may reflect which adversaries happened to train well as much as
any property of the architectures.

Two results do hold throughout. \texttt{CC-Simple-GNN} and \texttt{LC-Duelling-GNN} were
not harmed by any attack in any setting, for different reasons: the first barely changes
its routing under perturbation at all, whereas the second changes it and delivers the
packets regardless (Section~\ref{subsec:results_ind_flips}). GNN encoding is not
protective in itself, however, since the third GNN variant,
\texttt{CC-Duelling-GNN}, suffered the largest gap of the full-clamp independent attack.

\subsection{Further observations}
\label{subsec:conclusions_other}

One further result holds across the evaluation. The fraction of routing decisions an
attack changes does not predict the damage it does: attacks that changed nearly the same
fraction of decisions produced gaps several percentage points apart, so flip rates should
not be reported as attack success.

% ─────────────────────────────────────────────────────────────────────────────
\section{Limitations}
\label{sec:limitations}

\subsection{Statistical relevance of the results}
\label{subsec:limitations_statistics}

The most important limitation concerns how far the results generalise beyond the
particular adversaries that were trained. The learned adversary is not deterministic.
Its actor and critic are initialised at random, exploration noise is added to its
perturbations during training, and its updates are computed on randomly sampled
minibatches, and training was not seeded, so two training runs with the same
configuration produce different adversaries. Each of the 28 configurations of clamp,
scope and victim was trained only once, and every result in Chapter~\ref{chap:results}
describes that single trained adversary.

The confidence intervals reported there are computed over 30 paired evaluation episodes,
so they capture the variation between episodes for a given trained adversary, but not
the variation between training runs. A significant gap therefore shows that the adversary
that was trained beats its random control on these episodes, not that the attack
configuration would do so reliably if it were trained again. The results most exposed to
this are those that compare separately trained adversaries: the reversal on
\texttt{LC-Duelling}, where the independent parity-clamp adversary improved delivery, the
direction of the coordination effect, which compares the independent and coordinated
adversaries, the effect of the domain clamp, which compares adversaries trained under
each clamp, and the comparison with FGSM. The comparisons across timing budgets are less
exposed, since every budget, and the ungated run, evaluates the same trained adversary
under a different gate.

The size of the effects compounds this. Almost everything reported in
Chapter~\ref{chap:results} is a difference of one or two percentage points on a clean
delivery ratio of about $77\%$: no gated adversarial gap reaches $1.5$~pp, and the
coordination and clamp effects are of the same order. Perturbation alone, with no
adversarial direction, moves delivery by a comparable amount and in either direction,
which is why the random control is the reference rather than the clean arm. A paired
interval can resolve differences this small between two arms on the same episodes, and
the intervals reported are valid in that narrow sense, but resolving a difference is not
the same as attributing it: with one training run per configuration, a gap of this size
is as consistent with the idiosyncrasy of a single trained adversary as with a property
of the victim it was trained against. The derived quantities are more exposed still,
since the damage per attacked step divides one small quantity by another. The safe
reading of Chapter~\ref{chap:results} is that every gated attack left delivery within
about a percentage point of where a random perturbation of the same size left it, that
only the ungated attacks opened a wider margin, at most the $5.90$~pp on
\texttt{LC-Simple}, and that the differences \emph{between} configurations are not
established.

Two more statistical caveats apply. All evaluations use a single traffic seed, and
therefore the same 30 episodes, so the variation between traffic patterns is sampled only
by those episodes. And the chapter reports well over a hundred confidence intervals, each
at the $95\%$ level and without correction for multiple comparisons, so a few cells can be
expected to exclude zero by chance. Isolated significant cells, such as a single budget at
which an otherwise null variant separates from zero, should be read with caution, and the
conclusions of Section~\ref{sec:conclusions} rest on patterns that hold across budgets
rather than on individual cells.

\subsection{Scope of the evaluation}
\label{subsec:limitations_scope}

The evaluation covers a single topology, a single stressed traffic regime, one
perturbation budget of $\varepsilon = 0.30$ with
every agent compromised. The critical-state score was motivated partly by the thinning of
redundancy after link failures, yet no failures were injected in any evaluation, so that
part of its motivation is untested. Each adversary was
trained with the timing gate at $b = 0.25$ only, so the tighter budgets act on a subset of
its training states, and the ungated evaluation asks it to act on states it never
encountered. The reasons why the timing gate fired were recorded only for the FGSM-parity
runs.

\subsection{Design of the attack and its control}
\label{subsec:limitations_design}

The adversary is rewarded with the fraction of forwarding events that end in a drop, a
per-hop proxy for the end-to-end delivery on which it is scored, and the two can come
apart. Inspection of the trained adversaries also indicates that their outputs sit close
to the bounds of $[-1,1]$ for most features, so the attack perturbs most features by
nearly the full $\varepsilon$, whereas the random control draws each feature's
perturbation uniformly from $[-\varepsilon, \varepsilon]$. The adversarial gap therefore
reflects the size of the perturbation as well as its direction, and a random control
that also perturbs every feature by $\pm\varepsilon$, as the FGSM baseline's did, would
isolate direction more cleanly. Finally, the threat model assumes that the attacker can
train against a faithful simulator of the network and, when the attack is timed, read
the utilisation of every link while it attacks.

% ─────────────────────────────────────────────────────────────────────────────
\section{Future Work}
\label{sec:future_work}

\subsection{Repeated training for statistical relevance}
\label{subsec:future_repetition}

The limitation of Section~\ref{subsec:limitations_statistics} can be addressed directly.
The evaluation should be repeated with several independently trained adversaries per
configuration, each trained from a different random seed, so that the spread between
training runs can be measured and a result can be attributed to the attack rather than
to one run. That spread is what the present evaluation lacks, and measuring it is what
would show whether the differences of one or two percentage points reported here stand
above it at all.

\subsection{Optimisation-based PGD and SAJA attacks}
\label{subsec:future_optimisation}

This thesis took PGD, the strategically-timed attack of Lin et
al.~\cite{lin2019tacticsadversarialattackdeep} and SAJA as inspiration, but realised
both mechanisms through a
learned adversary rather than by solving an optimisation problem at each step. A natural
next step is to implement the attacks in their original form, as gradient optimisation
problems, and to evaluate them under the same protocol. A PGD
attack~\cite{madry2019deeplearningmodelsresistant} would solve for each attacked step's
perturbation by iterated projected gradient ascent on an attack objective, and could be
combined with the same timing gate. A SAJA-style attack~\cite{guo2025sajastateactionjointattack}
would perturb all compromised agents' observations jointly by gradient ascent, scoring
each candidate by the centralised critic's value of the joint action it induces.

Comparing these attacks with the learned adversary and with FGSM, within the same
admissible set and on the same paired episodes, would show whether learning the
adversary is beneficial for adversarial attacks on this controller, or whether direct
optimisation achieves more damage. Such a study would have to address several obstacles.
The routing decision is an $\operatorname{arg\,max}$ over candidate paths, so its gradient vanishes
almost everywhere and must be approximated, for instance with a straight-through
estimator. The local-critic variants provide no joint critic with which to score a
coordinated perturbation, so a different joint objective would be needed for them. And
the GNN encoder couples the agents' observations, so the gradient for one agent depends
on the others'. It would also have to account for the difference in cost, since an
optimisation-based attack requires many gradient evaluations of the victim at every
attacked step, whereas the learned adversary needs only a single forward pass once
trained and never needs the victim's gradients.
